\documentclass[a4paper,11pt]{article}
\usepackage[utf8]{inputenc}
\usepackage[T1]{fontenc}
\usepackage[french]{babel}
\usepackage{lmodern}
\usepackage{amsmath,amssymb,amsthm}
\usepackage{mathrsfs}
\usepackage{enumitem}
\usepackage{multicol}

\setlist[enumerate]{label=\textbf{\textcolor{Couleur8}{\arabic*.}}, left=1em}

% Si vous voulez aussi modifier les listes enumerate imbriquées :
\setlist[enumerate,2]{label=\textcolor{Couleur8}{\alph*)}}
\setlist[enumerate,3]{label=\textcolor{Couleur8}{\roman*)}}
\setlist[enumerate,4]{label=\textcolor{Couleur8}{\Alph*)}}
\usepackage{graphicx}
\usepackage{xcolor}
\usepackage{tcolorbox}
\usepackage{geometry}
\usepackage{fancyhdr}
\usepackage{titlesec}
\usepackage{cancel}
\usepackage{ifthen}
\usepackage{tikz}
\usetikzlibrary{arrows.meta}
\usepackage{tkz-tab}
\usepackage{eurosym}

% OPTION PROF/ÉLÈVE
% Décommentez UNE des deux lignes suivantes :
\newboolean{prof}\setboolean{prof}{true}   % Version professeur (avec corrections des devoirs)
%\newboolean{prof}\setboolean{prof}{false}  % Version élève (sans corrections des devoirs)

\definecolor{Couleur1}{HTML}{4A5D7A} %Th
\definecolor{Couleur2}{HTML}{A5B5D6} %Prop
\definecolor{Couleur3}{HTML}{A8C68A} %Méthode
\definecolor{Couleur6}{HTML}{8A9CC1} %Def
\definecolor{Couleur7}{HTML}{E6B459} %Rq Voc
\definecolor{Couleur8}{HTML}{D36740} %Titres
\definecolor{correction}{HTML}{D36740} % Couleur pour les corrections
\definecolor{coopmathscorr}{HTML}{F15929} % Couleur corrections coopmaths

% Configuration des marges
\geometry{
	top=2cm,
	bottom=2cm,
	inner=1.5cm,
	outer=1.5cm,
}

% Police
\renewcommand{\familydefault}{\sfdefault}

% Compteurs
\newcounter{seancenum}
\newcounter{seancenumD}
\setcounter{seancenum}{1}
\setcounter{seancenumD}{1}

% Configuration des en-têtes et pieds de page
\pagestyle{fancy}
\fancyhf{}
\ifthenelse{\boolean{prof}}{
    \fancyhead[L]{\textcolor{Couleur8}{\textbf{Corrections Automatismes - Classe de seconde - Version Professeur}}}
}{
    \fancyhead[L]{\textcolor{Couleur8}{\textbf{Corrections Automatismes - Séance 6 - Classe de seconde}}}
}
\fancyhead[R]{\textcolor{Couleur8}{\textbf{2025-2026}}}
\fancyfoot[L]{\textcolor{Couleur8}{Mme Bertail et Mme Pinto}}
\fancyfoot[R]{\textcolor{Couleur8}{\thepage}}
\renewcommand{\headrulewidth}{1pt}
\renewcommand{\headrule}{\hbox to\headwidth{\color{Couleur8}\leaders\hrule height \headrulewidth\hfill}}

% Environnements personnalisés
\newtcolorbox{seance}[1]{
    colback=Couleur6!10,
    colframe=Couleur1!65,
    fonttitle=\bfseries\large,
    title=Correction Séance \theseancenum #1,
    left=5mm,
    right=5mm,
    top=3mm,
    bottom=3mm,
    arc=0mm,
    after=\stepcounter{seancenum}
}

\newtcolorbox{seanceD}[1]{
    colback=Couleur6!5,
    colframe=Couleur6!45,
    fonttitle=\bfseries\large,
    title=Correction Devoirs - Séance \theseancenumD #1,
    left=5mm,
    right=5mm,
    top=3mm,
    bottom=3mm,
    arc=0mm,
    after=\stepcounter{seancenumD}
}

\begin{document}

\setcounter{seancenum}{6}
\newpage
\begin{seance}{} %6

\begin{enumerate}
\item Le nombre de solutions de l'équation $f(x)=0$ est le nombre d'antécédents de $0$ par la fonction $f$.\\
Puisque la droite d'équation $y = 0$ (droite horizontale) coupe $1$ fois la courbe, on en déduit que l'équation $f(x)=0$ admet ${\color{correction}\boldsymbol{1}}$ {\bfseries \color{correction}solution}.

\item On a :\\
          $f(-3)=1- 2\times (-3)^2={\color{correction}\boldsymbol{-17}}$\\{\color{Couleur1} Mentalement : \\
    On commence par "calculer" le carré de $-3$ :  $(-3)^2=9$.\\
    Puis on multiplie le résultat par $2$ : $2\times 9=18$.\\
    On calcule alors : $1-18=-17$.}

\item On remplace $x$ par $x+3$ dans l'expression de $u$ : \\
$u(x+3)=-3(x+3)-10=-3x-9-10={\color{correction}\boldsymbol{-3x-19}}$

\item $(-1)^{n+5}=(-1)^{4}\times (-1)^{n+1}=1\times (-1)^{n+1} = {\color{correction}\boldsymbol{(-1)^{n+1}}}$

La bonne réponse est la réponse {\bfseries \color{correction}C}.

\item $N=\dfrac{15^7}{5^4}=\dfrac{5^7\times 3^7 }{5^4}=3^7\times 5^{3}=3^4\times3^3\times5^3=3^4\times(3\times5)^3=3^{4}\times15^3 ={\color{correction}\boldsymbol{81\times 15^{3}}}$

La bonne réponse est la réponse {\bfseries \color{correction}B}.

\item $A = \dfrac{1}{1-\dfrac{3}{5}} = \dfrac{1}{\dfrac{5}{5}-\dfrac{3}{5}} = \dfrac{1}{\dfrac{2}{5}} = 1 \times \dfrac{5}{2} = {\color{correction}\boldsymbol{\dfrac{5}{2}}}$
\end{enumerate}

\end{seance}

\end{document}